% GENERATED FILE. Edit canonical lessons, then run npm run generate:books.
% canonical-source-hash: fnv1a64:124baad0d19d766f
% canonical-lessons: SA-C356-vijnapayati, SA-C356-niyojayati, SA-C356-anumodate, SA-C356-vibhagah, SA-C356-prastavah

\chapter{Notices from an Office}
\label{ch:sa-notices-from-an-office}
\hlchaptermodality{356}

\begin{chapteropening}
\textbf{By the end of this chapter:} \emph{I can say what an office made known and whom it appointed.}

Complete the last lesson of chapter 356: I can say what an office made known and whom it appointed.
\end{chapteropening}

\section[\texorpdfstring{\sk{विज्ञापयति} (vijñāpayati)}{विज्ञापयति (vijñāpayati)}]{\sk{विज्ञापयति} (vijñāpayati) — notifies; makes a formal request}

\label{lesson:SA-C356-vijnapayati}



\begin{quote}
Before the new one: say the Sanskrit for a president, then the Sanskrit for consults.
\end{quote}

\subsection*{You'll want to know: \sk{विज्ञापयति}}
\textbf{\sk{विज्ञापयति}} — \emph{vijñāpayati} — \textquotedblleft{}notifies; makes a formal request\textquotedblright{}.

\begin{quote}
\sk{कार्यालयः} \sk{नूतनं} \sk{नियमं} \sk{विज्ञापयति।}
\end{quote}

\begin{quote}
\emph{kāryālayaḥ nūtanaṃ niyamaṃ vijñāpayati.}
\end{quote}

\begin{quote}
The office gives notice of a new rule.
\end{quote}

Past (he …): \sk{सः} \sk{अध्यक्षं} \sk{विज्ञापितवान्}, \emph{saḥ adhyakṣaṃ vijñāpitavān}.

\begin{cousinweb}
\textbf{\sk{विज्ञापयति}} (\emph{vijñāpayati}) is built on \emph{vi} and the root √jñā, \textquotedblleft{}to know\textquotedblright{}, in its causative form, \textquotedblleft{}to make known\textquotedblright{}; MW gives the senses \textquotedblleft{}to inform\textquotedblright{}, \textquotedblleft{}to apprise\textquotedblright{} and \textquotedblleft{}to request\textquotedblright{} (MW).
\end{cousinweb}

\begin{grammarlens}[title={What you've built}]
One more word for giving reasons and opinions.
\end{grammarlens}

\subsection*{Your turn}
\begin{itemize}
\raggedright
  \item \emph{Say it:} \emph{vijñāpayati}
  \item \emph{Say it:} \emph{vijñāpayati}, once more
  \item \emph{Say it:} \emph{kāryālayaḥ nūtanaṃ niyamaṃ vijñāpayati.} — The office gives notice of a new rule.
  \item \emph{Recall:} say the Sanskrit for a president, then the Sanskrit for consults, then say \emph{vijñāpayati} again
\end{itemize}

\begin{tcolorbox}[breakable,colback=teal!4,colframe=teal!35!black,title={Before you move on}]
What does \sk{विज्ञापयति} mean? (\textquotedblleft{}Notifies; makes a formal request\textquotedblright{}.) Say it once more.
\end{tcolorbox}
\section[\texorpdfstring{\sk{नियोजयति} (niyojayati)}{नियोजयति (niyojayati)}]{\sk{नियोजयति} (niyojayati) — appoints}

\label{lesson:SA-C356-niyojayati}



\begin{quote}
Before the new one: say the Sanskrit for consults, then the Sanskrit for notifies.
\end{quote}

\subsection*{You'll want to know: \sk{नियोजयति}}
\textbf{\sk{नियोजयति}} — \emph{niyojayati} — \textquotedblleft{}appoints\textquotedblright{}.

\begin{quote}
\sk{राजा} \sk{नूतनं} \sk{लिपिकं} \sk{नियोजयति।}
\end{quote}

\begin{quote}
\emph{rājā nūtanaṃ lipikaṃ niyojayati.}
\end{quote}

\begin{quote}
The king appoints a new clerk.
\end{quote}

Past (he …): \sk{सः} \sk{सदस्यं} \sk{नियोजितवान्}, \emph{saḥ sadasyaṃ niyojitavān}.

\begin{cousinweb}
\textbf{\sk{नियोजयति}} (\emph{niyojayati}) is built on \emph{ni}, \textquotedblleft{}down, in\textquotedblright{}, and the root √yuj, \textquotedblleft{}to join\textquotedblright{}, in its causative form; MW gives \emph{ni-√yuj} the senses \textquotedblleft{}to enjoin, to order\textquotedblright{} and \textquotedblleft{}to appoint, to instal\textquotedblright{} (MW).
\end{cousinweb}

\begin{grammarlens}[title={What you've built}]
One more word for giving reasons and opinions.
\end{grammarlens}

\subsection*{Your turn}
\begin{itemize}
\raggedright
  \item \emph{Say it:} \emph{niyojayati}
  \item \emph{Say it:} \emph{niyojayati}, once more
  \item \emph{Say it:} \emph{rājā nūtanaṃ lipikaṃ niyojayati.} — The king appoints a new clerk.
  \item \emph{Recall:} say the Sanskrit for consults, then the Sanskrit for notifies, then say \emph{niyojayati} again
\end{itemize}

\begin{tcolorbox}[breakable,colback=teal!4,colframe=teal!35!black,title={Before you move on}]
What does \sk{नियोजयति} mean? (\textquotedblleft{}Appoints\textquotedblright{}.) Say it once more.
\end{tcolorbox}
\section[\texorpdfstring{\sk{अनुमोदते} (anumodate)}{अनुमोदते (anumodate)}]{\sk{अनुमोदते} (anumodate) — approves of}

\label{lesson:SA-C356-anumodate}



\begin{quote}
Before the new one: say the Sanskrit for notifies, then the Sanskrit for appoints.
\end{quote}

\subsection*{You'll want to know: \sk{अनुमोदते}}
\textbf{\sk{अनुमोदते}} — \emph{anumodate} — \textquotedblleft{}approves of\textquotedblright{}.

\begin{quote}
\sk{अध्यक्षः} \sk{योजनाम्} \sk{अनुमोदते।}
\end{quote}

\begin{quote}
\emph{adhyakṣaḥ yojanām anumodate.}
\end{quote}

\begin{quote}
The president approves the plan.
\end{quote}

Past (he …): \sk{सः} \sk{योजनाम्} \sk{अनुमोदितवान्}, \emph{saḥ yojanām anumoditavān}.

\begin{cousinweb}
\textbf{\sk{अनुमोदते}} (\emph{anumodate}) is built on \emph{anu}, \textquotedblleft{}after, along\textquotedblright{}, and the root √mud, \textquotedblleft{}to rejoice\textquotedblright{}; MW gives \emph{anu-√mud} \textquotedblleft{}to join in rejoicing\textquotedblright{} and \textquotedblleft{}to express approval, to applaud\textquotedblright{} (MW).
\end{cousinweb}

\begin{grammarlens}[title={What you've built}]
One more word for giving reasons and opinions.
\end{grammarlens}

\subsection*{Your turn}
\begin{itemize}
\raggedright
  \item \emph{Say it:} \emph{anumodate}
  \item \emph{Say it:} \emph{anumodate}, once more
  \item \emph{Say it:} \emph{adhyakṣaḥ yojanām anumodate.} — The president approves the plan.
  \item \emph{Recall:} say the Sanskrit for notifies, then the Sanskrit for appoints, then say \emph{anumodate} again
\end{itemize}

\begin{tcolorbox}[breakable,colback=teal!4,colframe=teal!35!black,title={Before you move on}]
What does \sk{अनुमोदते} mean? (\textquotedblleft{}Approves of\textquotedblright{}.) Say it once more.
\end{tcolorbox}
\section[\texorpdfstring{\sk{विभागः} (vibhāgaḥ)}{विभागः (vibhāgaḥ)}]{\sk{विभागः} (vibhāgaḥ) — a department; a division}

\label{lesson:SA-C356-vibhagah}



\begin{quote}
Before the new one: say the Sanskrit for appoints, then the Sanskrit for approves of.
\end{quote}

\subsection*{You'll want to know: \sk{विभागः}}
\textbf{\sk{विभागः}} — \emph{vibhāgaḥ} — \textquotedblleft{}a department; a division\textquotedblright{}.

\begin{quote}
\sk{कार्यालये} \sk{पञ्च} \sk{विभागाः} \sk{सन्ति।}
\end{quote}

\begin{quote}
\emph{kāryālaye pañca vibhāgāḥ santi.}
\end{quote}

\begin{quote}
The office has five departments.
\end{quote}

A masculine noun.

\begin{cousinweb}
\textbf{\sk{विभागः}} (\emph{vibhāgaḥ}) is built on \emph{vi}, \textquotedblleft{}apart\textquotedblright{}, and the root √bhaj, \textquotedblleft{}to divide, to share\textquotedblright{}; MW gives \emph{vi-bhāga} \textquotedblleft{}a share, a section, a division\textquotedblright{}, and \textquotedblleft{}a department\textquotedblright{} is its modern sense (MW).
\end{cousinweb}

\begin{grammarlens}[title={What you've built}]
One more word for giving reasons and opinions.
\end{grammarlens}

\subsection*{Your turn}
\begin{itemize}
\raggedright
  \item \emph{Say it:} \emph{vibhāgaḥ}
  \item \emph{Say it:} \emph{vibhāgaḥ}, once more
  \item \emph{Say it:} \emph{kāryālaye pañca vibhāgāḥ santi.} — The office has five departments.
  \item \emph{Recall:} say the Sanskrit for appoints, then the Sanskrit for approves of, then say \emph{vibhāgaḥ} again
\end{itemize}

\begin{tcolorbox}[breakable,colback=teal!4,colframe=teal!35!black,title={Before you move on}]
What does \sk{विभागः} mean? (\textquotedblleft{}A department; a division\textquotedblright{}.) Say it once more.
\end{tcolorbox}
\section[\texorpdfstring{\sk{प्रस्तावः} (prastāvaḥ)}{प्रस्तावः (prastāvaḥ)}]{\sk{प्रस्तावः} (prastāvaḥ) — a proposal}

\label{lesson:SA-C356-prastavah}



\begin{quote}
Before the new one: say the Sanskrit for approves of, then the Sanskrit for a department.
\end{quote}

\subsection*{You'll want to know: \sk{प्रस्तावः}}
\textbf{\sk{प्रस्तावः}} — \emph{prastāvaḥ} — \textquotedblleft{}a proposal\textquotedblright{}.

\begin{quote}
\sk{अध्यक्षः} \sk{प्रस्तावं} \sk{स्वीकरोति।}
\end{quote}

\begin{quote}
\emph{adhyakṣaḥ prastāvaṃ svīkaroti.}
\end{quote}

\begin{quote}
The president accepts the proposal.
\end{quote}

A masculine noun.

\begin{cousinweb}
\textbf{\sk{प्रस्तावः}} (\emph{prastāvaḥ}) is built on \emph{pra} and the root √stu, \textquotedblleft{}to praise\textquotedblright{}; MW gives \emph{pra-stāva} \textquotedblleft{}an introduction\textquotedblright{} and \textquotedblleft{}introducing a topic\textquotedblright{}, and \textquotedblleft{}a proposal\textquotedblright{} is its modern sense (MW).
\end{cousinweb}

\begin{grammarlens}[title={What you've built}]
One more word for giving reasons and opinions.
\end{grammarlens}

\subsection*{Your turn}
\begin{itemize}
\raggedright
  \item \emph{Say it:} \emph{prastāvaḥ}
  \item \emph{Say it:} \emph{prastāvaḥ}, once more
  \item \emph{Say it:} \emph{adhyakṣaḥ prastāvaṃ svīkaroti.} — The president accepts the proposal.
  \item \emph{Recall:} say the Sanskrit for approves of, then the Sanskrit for a department, then say \emph{prastāvaḥ} again
\end{itemize}

\begin{tcolorbox}[breakable,colback=teal!4,colframe=teal!35!black,title={Before you move on}]
What does \sk{प्रस्तावः} mean? (\textquotedblleft{}A proposal\textquotedblright{}.) Say it once more.
\end{tcolorbox}
